\section{Introduction}
\label{sec:introduction}
Automated warehouses mix workers, forklifts, pallet trucks, and autonomous mobile robots (AMRs) in shared aisles, so multi-class object detection is a basic requirement for their safe operation. Detectors in this setting face large scale and viewpoint variation, repetitive racking with frequent occlusion, and the compute, memory, and thermal limits of single-board computers such as the Raspberry Pi~4.

Lightweight convolutional detectors such as YOLO11s~\cite{ultralytics2024yolo11} run efficiently on such devices, but each convolution aggregates only local context. YOLO11s therefore adds one self-attention block, \texttt{C2PSA}, at its lowest-resolution stage $P_5$; it holds $990{,}976$ parameters (about $10.5\%$ of the model), although at the $20 \times 20$ resolution of $P_5$ ($N = 400$ tokens) its quadratic cost $\mathcal{O}(N^2)$ is small in absolute terms. Selective state space models (SSMs) such as Mamba~\cite{gu2023mamba} and their visual variants~\cite{liu2024vmamba,zhu2024visionmamba} offer global context with cost linear in $N$, and several works have inserted them into YOLO-style detectors~\cite{wang2025mambayolo,lei2025mambanext,chan2026mambapsa}. These works report accuracy, FLOPs, or GPU and Jetson throughput; the measured cost of a plugin-free scan on an embedded CPU, and what determines it, has received little attention.

We therefore ask three research questions. \textbf{RQ1:} Does replacing the $P_5$ attention block of YOLO11s with a simplified, plugin-free 2D selective-scan block preserve accuracy while reducing parameters and FLOPs? \textbf{RQ2:} Does inserting a newly initialized scan block into a pretrained detector require a dedicated protocol (two-stage training, identity initialization)? \textbf{RQ3:} What latency and memory does a plugin-free selective scan incur on an embedded CPU, and which implementation factors determine that cost?

This is primarily a systems and measurement study of deployment cost; the block itself is a simplified variant of existing designs. Our contributions are as follows:
\begin{itemize}
    \item \textbf{A hardware-counter autopsy of embedded SSM cost} on a Raspberry Pi~4 (ARM Cortex-A72). Despite reducing theoretical complexity ($-6.72\%$ parameters, $-0.44$ GFLOPs), the SS2D detector is $18\%$ slower than the attention baseline ($1314$ vs.\ $1117$\,ms with four threads) and uses $60\%$ more peak memory ($373.6$ vs.\ $233.4$\,MiB). Using controlled graph interventions and hardware counters (\texttt{perf stat}), we prove that this penalty does not stem from arithmetic FLOPs, but from operator dispatch overhead in the unrolled sequential scan ($14{,}771$ nodes) and severe memory-bandwidth pressure in the parallel scan ($2.2\times$ DRAM traffic, $2.5\times$ L2 refills).
    \item \textbf{A fine-grained per-component cost breakdown of replacing \texttt{C2PSA}} with a plugin-free 4-direction selective scan. We reveal that the $63.9\%$ Layer-10 parameter saving stems from narrower channel projections rather than the token-mixing core, and that redundant per-direction parameter projection makes the core more expensive than self-attention at $L = 400$.
    \item \textbf{An analysis of CNN--SSM grafting pitfalls.} We show that zero-initializing a projection convolution ahead of batch normalization preserves identity only at step 0, explaining the early-epoch gradient instability that necessitated two-stage warmup, and show that the trained scan acts primarily as a localized mixer.
\end{itemize}
